\documentclass[a4paper]{article}
\usepackage[pdftex]{graphicx}
\usepackage[utf8]{inputenc}
\usepackage{enumerate}
\usepackage{colortbl}
\usepackage{href-ul}
\hypersetup{
	colorlinks=true,
	linkcolor=blue,
	urlcolor=blue}
\usepackage{siunitx}
\sisetup{locale=DE}
\usepackage{icomma}
\usepackage{amssymb}
\usepackage{geometry}
\geometry{a4paper, top=15mm, left=15mm, right=15mm, bottom=15mm,
	headsep=10mm, footskip=12mm}

\begin{document}
	{\bf Kinetic Energy}
	
	\begin{enumerate}[1.]
		{\bf \item The following bodies with characteristic masses and velocities are given. Assign them to the table.}:\\
		\fbox{Football}, \fbox{Shooting Star}, \fbox{Car}, \fbox{ICE}, \fbox{Table Tennis Ball}, \fbox{Fighter Jet}, \fbox{Raindrop}, \fbox{Sprinter}
		
		\renewcommand{\arraystretch}{2}
		\begin{tabular}{|>{\columncolor[gray]{.8}}p{3 cm}|p{3 cm}|p{3cm}|p{3 cm}|}
			\hline
			\rowcolor[gray]{.8} Body & Mass &Velocity & Kinetic Energy \\
			\hline
			&$ \SI{0.05}{\gram} $&$ \SI{5}{\meter\per\second} $&\\
			\hline
			&$ \SI{4,7}{\gram} $&$ \SI{10}{\meter\per\second} $& \\
			\hline
			&$ \SI{75}{\kilo\gram} $&$ \SI{10}{\meter\per\second} $&\\
			\hline
			&$ \SI{0,1}{\gram} $&$ \SI{30}{\kilo\meter\per\second} $& \\
			\hline
			&$ \SI{0.42}{\kilo\gram} $&$ \SI{100}{\kilo\meter\per\hour} $& \\
			\hline
			&$ \SI{1}{\tonne} $&$ \SI{130}{\kilo\meter\per\hour} $& \\
			\hline
			&$ \SI{2}{\tonne} $&$ \SI{1200}{\kilo\meter\per\hour} $&\\
			\hline
			&$ \SI{300}{\tonne} $&$ \SI{250}{\kilo\meter\per\hour} $&\\
			\hline
		\end{tabular}
		
		
		
		{\bf \item A roller coaster car with mass $m=\SI{200}{\kilogram}$ goes through a loop.} To do this, it starts with the speed $v_0 = 0$ at point A at a height of $h= \SI{30}{\meter}$ (See sketch, not to scale).
		
		\begin{enumerate}[a)]
			\item Calculate the speed of the car at points B and D (ignoring friction)
			\item How would these speeds change if the car had more than $\SI{200}{\kilogram}$ mass (ignoring friction)?
			\item The car should be brought to a standstill behind D over a distance of $s=\SI{25}{\meter}$. What is the required braking force?
		\end{enumerate}
		
		\includegraphics[width=9 cm]{achterbahn115.pdf}
		
		
		{\bf \item task\\}
		\begin{enumerate}[a)]
			\item Calculate the speed at which a high diver jumps from the $\SI{10}{\meter}$ tower
			hits the water.
			\item From what height did someone jump who reached a speed of $v=\SI{10}{\meter\per\second}$ shortly before entering the water?
		\end{enumerate}
	\end{enumerate}
	
	\fbox{
		\begin{minipage}{0.5\textwidth}
			For the solution, please \href{https://www.okuyakl.de/phys/p8ekiL411/le411.pdf}{click here} or scan the QR code.\\
			Additional worksheets are available at
			
			\href{http://www.okuyakl.com}{www.okuyakl.com}
		\end{minipage}
		\hfill
		\begin{minipage}{0.4\textwidth}
			\includegraphics[width=1.5 cm]{../../viecher/zwe03}
			\includegraphics[width=3 cm]{qre411}
			\includegraphics[width=2 cm]{../../viecher/afanticon1}
			
	\end{minipage}}
\end{document}